\section{Resultados esperados}

Se espera que YOLOv11 en INT8 ofrezca el mejor compromiso: mayor mAP@0.5 y F1
macro que YOLOv8n, ventaja que se conservaría en condiciones adversas, con menor
latencia, memoria y tamaño. La cuantización INT8 debería reducir latencia y
memoria con una degradación de mAP no superior al umbral que se fije en la
corrida piloto. En la fase pedagógica se espera que la revisión asistida reduzca
el tiempo por página y la carga mental.

\begin{table}[H]
  \centering
  \caption{Resultados esperados por configuración}
  \label{tab:resultados}
  \small
  \begin{tabularx}{\columnwidth}{@{}lX@{}}
    \toprule
    \textbf{Configuración} & \textbf{Resultado esperado} \\
    \midrule
    YOLOv8n FP32 & Línea base de precisión y latencia \\
    YOLOv8n INT8 & Menor latencia y memoria \\
    YOLOv11 FP32 & Mayor precisión que YOLOv8n \\
    YOLOv11 INT8 & Mejor compromiso precisión--eficiencia \\
    \bottomrule
  \end{tabularx}
  \notaTabla{Elaboración propia.}
\end{table}
